\chapter{The Chinese Writing System and Its Spread}

\section{Why writing comes first}

The single most important shared element of East Asian civilization is the Chinese
script. Until the late nineteenth century an educated person in Seoul, Kyoto or Hanoi could
read the same Classical Chinese texts and communicate in writing (``brush talk'') with a
Chinese scholar without sharing a spoken language. Writing carried everything else in this
course: Confucian classics, Buddhist sutras, law codes and the examination system.

\begin{definition}[Literary (Classical) Chinese]
\emph{Literary Chinese} (\zht{文言文}, wényánwén) is the terse written language based on
the texts of the late Zhou and Han periods. It was never anyone's spoken language after
antiquity, and it remained the prestige written language of the whole region until the
twentieth century, much as Latin did in medieval Europe.
\end{definition}

\section{Origins}

\begin{itemize}
  \item \textbf{Oracle bones} (\zht{甲骨文}, jiǎgǔwén), late Shang dynasty (c.\ 1250--1050
        BCE): questions to ancestors were carved on turtle plastrons and ox scapulae, the
        bone was heated, and the cracks were read as answers. This is the earliest
        substantial body of Chinese writing.
  \item \textbf{Bronze inscriptions} (Zhou): records cast inside ritual vessels.
  \item \textbf{Qin standardization} (221 BCE): the First Emperor imposed a single
        ``small seal'' script, along with standard weights, measures and axle widths.
  \item \textbf{Clerical and standard scripts} (Han onward): the brush-written forms that
        are recognizably modern.
\end{itemize}

\begin{important}[Characters are not ``pictures of ideas'']
A common misconception is that Chinese characters are pictographs or ``ideograms''. Only a
small fraction began as pictures, and over 80\% of characters are \emph{phono-semantic
compounds}: one component gives a hint to meaning, another a hint to pronunciation. The
script is \emph{logographic}: each character writes a \emph{morpheme} (a meaningful
syllable) of the language.
\end{important}

\section{How characters are built}

Traditional analysis (from the Han dictionary \emph{Shuowen jiezi}, c.\ 100 CE) recognized
six types. We only need four.

\begin{center}
\begin{tabular}{llll}
\toprule
Type & Example & Pinyin & How it works \\
\midrule
Pictograph & \zht{山} \quad \zht{木} & shān, mù & stylized drawing of a mountain, a tree \\
Simple indicative & \zht{上} \quad \zht{本} & shàng, běn & a mark above a line; a mark at the root of a tree \\
Compound ideograph & \zht{休} \quad \zht{明} & xiū, míng & person + tree = rest; sun + moon = bright \\
Phono-semantic & \zht{媽} \quad \zht{河} & mā, hé & semantic radical + sound component \\
\bottomrule
\end{tabular}
\end{center}

\subsection{A phono-semantic family}

The component \zht{青} (qīng, ``blue-green'') gives the sound to a whole family of
characters; the radical on the left tells you the field of meaning.

\begin{center}
\begin{tikzpicture}[comp/.style={draw, thick, minimum size=0.9cm, font=\large},
                    lab/.style={font=\scriptsize, align=center}]
  \node[comp, fill=yellow!20] (q) at (0,0) {\zht{青}};
  \node[lab, below=0.05cm of q] {qīng\\sound component};
  \foreach \rad/\full/\py/\gl/\x/\i in {%
      \zht{氵}/\zht{清}/qīng/clear (water)/-4.5/1,
      \zht{日}/\zht{晴}/qíng/sunny (sun)/-1.5/2,
      \zht{目}/\zht{睛}/jīng/eyeball (eye)/1.5/3,
      \zht{言}/\zht{請}/qǐng/to request (speech)/4.5/4} {
    \node[comp, fill=blue!10] (r\i) at (\x-0.5,2.4) {\rad};
    \node[comp, fill=yellow!20] (s\i) at (\x+0.5,2.4) {\zht{青}};
    \node[comp, fill=green!12, minimum size=1.1cm] (f\i) at (\x,4.2) {\full};
    \node[lab, above=0.05cm of f\i] {\py\\\gl};
    \draw[->] (r\i.north) -- (f\i.south west);
    \draw[->] (s\i.north) -- (f\i.south east);
    \draw[gray, dashed] (q.north) -- (s\i.south);
  }
\end{tikzpicture}
\end{center}

\begin{remark}
Sound changes over 2000 years mean that the phonetic hint is often only approximate (as
with jīng above). In Korean and Japanese readings the hint sometimes works better than in
modern Mandarin, because those readings preserve older pronunciations.
\end{remark}

\subsection{Radicals}

Dictionaries are organized by \emph{radicals} (\zht{部首}, bùshǒu), the semantic
components. The Kangxi dictionary (1716) fixed the standard list of 214 radicals that is
still used in Japanese and Korean character dictionaries.

\begin{center}
\begin{tabular}{llll}
\toprule
Radical & Form as component & Meaning & Example \\
\midrule
\zht{水} & \zht{氵} & water & \zht{海} sea, \zht{洋} ocean \\
\zht{人} & \zht{亻} & person & \zht{仁} benevolence, \zht{信} trust \\
\zht{心} & \zht{忄} & heart/mind & \zht{情} feeling, \zht{性} nature \\
\zht{言} & \zht{訁} & speech & \zht{語} language, \zht{說} speak \\
\zht{木} & \zht{木} & tree/wood & \zht{林} grove, \zht{森} forest \\
\bottomrule
\end{tabular}
\end{center}

\section{Simplified and traditional characters}

The People's Republic introduced \emph{simplified characters} in 1956 and 1964 to raise
literacy. Taiwan, Hong Kong and Macau kept \emph{traditional} forms. Japan simplified a
smaller set independently after 1946 (\ja{新字体}, shinjitai), and South Korea kept
traditional forms for the hanja it still uses.

\begin{center}
\begin{tabular}{lllll}
\toprule
Meaning & Traditional & Simplified (PRC) & Japanese & Korean hanja \\
\midrule
country & \zht{國} & \zh{国} & \ja{国} & \zht{國} \\
learn & \zht{學} & \zh{学} & \ja{学} & \zht{學} \\
dragon & \zht{龍} & \zh{龙} & \ja{竜} & \zht{龍} \\
east & \zht{東} & \zh{东} & \ja{東} & \zht{東} \\
\bottomrule
\end{tabular}
\end{center}

\section{The script travels: Korea, Japan, Vietnam}

The basic problem for neighbours was that Chinese is an isolating language (no
inflection, word order does the grammar) while Korean and Japanese are agglutinative
(long strings of suffixes and particles, verb at the end). Writing a Japanese or Korean
sentence in Chinese characters was awkward. Each society developed solutions in roughly
the same three stages.

\begin{enumerate}
  \item \textbf{Write in Literary Chinese} and read it in Chinese order.
  \item \textbf{Annotate Chinese} with marks that show how to read it in native word order
        with native endings (Japanese \ja{漢文訓読}, kanbun kundoku; Korean \ko{구결},
        gugyeol).
  \item \textbf{Borrow characters for their sound} to write native words, then simplify
        those borrowed characters into a separate script (kana), or invent a new script
        altogether (hangul).
\end{enumerate}

\begin{center}
\begin{tikzpicture}[node distance=1.2cm and 1.6cm,
    st/.style={draw, rounded corners, thick, align=center, font=\small, minimum width=2.6cm},
    ar/.style={->, thick}]
  \node[st, fill=orange!20] (cn) {China\\Chinese characters\\(Han, 1st c.\ CE+)};
  \node[st, fill=green!15, right=of cn, yshift=1.6cm] (kr) {Korea\\idu, hyangchal\\(5th--7th c.)};
  \node[st, fill=red!12, right=of kr] (jp) {Japan\\man'yōgana\\(7th--8th c.)};
  \node[st, fill=yellow!20, right=of cn, yshift=-1.6cm] (vn) {Vietnam\\chữ Nôm\\(13th c.+)};
  \node[st, fill=green!30, below=of kr, yshift=0.4cm, xshift=1.2cm] (hg) {hangul\\(1443/1446)};
  \node[st, fill=red!25, right=of jp, xshift=-0.6cm] (kana) {hiragana /\\katakana\\(9th c.)};
  \draw[ar] (cn) -- (kr);
  \draw[ar] (kr) -- node[above, font=\scriptsize] {via scribes} (jp);
  \draw[ar] (cn) -- (vn);
  \draw[ar] (jp) -- (kana);
  \draw[ar, dashed] (kr) -- node[right, font=\scriptsize, align=left] {invented,\\not derived} (hg);
\end{tikzpicture}
\end{center}

\subsection{Japan: kanji and kana}

Japan received writing through Korean scribes (traditionally the scholar Wani from
Baekje, c.\ 400 CE). The \emph{Man'yōshū} poetry anthology (8th c.) used characters
purely for sound (\ja{万葉仮名}, man'yōgana). These were then simplified:

\begin{itemize}
  \item \textbf{Hiragana} (\ja{ひらがな}): cursive simplifications of whole characters,
        associated with court women's writing (e.g.\ \ja{安} $\to$ \ja{あ},
        \ja{以} $\to$ \ja{い}).
  \item \textbf{Katakana} (\ja{カタカナ}): fragments of characters used by monks to annotate
        Buddhist texts (e.g.\ \ja{阿} $\to$ \ja{ア}, \ja{伊} $\to$ \ja{イ}).
\end{itemize}

Modern Japanese mixes all three: kanji for content words, hiragana for grammatical
endings and particles, katakana for loanwords and emphasis. A character typically has an
\emph{on} reading (borrowed Chinese sound) and a \emph{kun} reading (native Japanese
word): \ja{山} is read \emph{san} in \ja{富士山} (Fujisan) but \emph{yama} on its own.

\begin{example}[One sentence, three scripts]
\ja{私はコーヒーを飲みます。} (Watashi wa kōhī o nomimasu, ``I drink coffee.'')
\ja{私} and \ja{飲} are kanji; \ja{コーヒー} is katakana for an English loanword;
\ja{は}, \ja{を}, \ja{みます} are hiragana particles and verb endings.
\end{example}

\subsection{Korea: hanja, idu and hangul}

Korean elites wrote Literary Chinese for over a thousand years. To record Korean they
used \emph{idu} (\ko{이두}), a system using characters for Korean grammatical elements, and
\emph{hyangchal} for native poetry (the \emph{hyangga} of Silla). In 1443 King Sejong and
scholars of the Hall of Worthies created a new alphabet, promulgated in 1446 as
\emph{Hunmin jeongeum} (\ko{훈민정음}, ``the correct sounds for instructing the
people''). We return to hangul in the Korea lecture; here note:

\begin{itemize}
  \item It is a true alphabet, but letters are grouped into syllable blocks so that a
        text \emph{looks} like a page of characters.
  \item Consonant shapes depict the speech organs (\ko{ㄱ} the tongue root blocking the
        throat, \ko{ㅁ} the mouth).
  \item Elites resisted it for centuries; it became the standard national script only in
        the twentieth century. South Korea now uses very few hanja; North Korea none.
\end{itemize}

\subsection{Vietnam}

Vietnam wrote Literary Chinese and later developed \emph{chữ Nôm}, new characters built
from Chinese components to write Vietnamese. Under French rule a Roman alphabet
(\emph{quốc ngữ}) devised by missionaries replaced both.

\section{Sino-xenic vocabulary}

Because all four societies borrowed huge numbers of words from Chinese, they share a
layer of ``Sino-'' vocabulary. Recognizing it is a good party trick and shows the depth of
shared heritage.

\begin{center}
\begin{tabular}{lllll}
\toprule
Characters & Mandarin & Japanese & Korean & Meaning \\
\midrule
\zht{大學} & dàxué & daigaku (\ja{大学}) & daehak (\ko{대학}) & university \\
\zht{學生} & xuéshēng & gakusei (\ja{学生}) & haksaeng (\ko{학생}) & student \\
\zht{經濟} & jīngjì & keizai (\ja{経済}) & gyeongje (\ko{경제}) & economy \\
\zht{文化} & wénhuà & bunka (\ja{文化}) & munhwa (\ko{문화}) & culture \\
\zht{國家} & guójiā & kokka (\ja{国家}) & gukka (\ko{국가}) & state \\
\bottomrule
\end{tabular}
\end{center}

\begin{aside}
Many of these words (\emph{economy}, \emph{society}, \emph{philosophy}, \emph{science}) were
actually coined in \emph{Meiji Japan} to translate Western concepts and then borrowed
\emph{back} into Chinese and Korean around 1900. The flow of the character sphere reversed
in the modern period. We will see this again in Week 9.
\end{aside}

\section{Writing, power and literacy}

\begin{itemize}
  \item Mastery of characters required years of study, so literacy was a marker of elite
        status (scholar-officials, \emph{yangban}, court nobles, monks).
  \item Calligraphy (\zht{書法} shūfǎ, \ja{書道} shodō, \ko{서예} seoye) was an art form and
        a test of moral cultivation.
  \item Vernacular scripts (kana, hangul) opened writing to women and commoners and
        produced major literatures (\emph{The Tale of Genji}; Korean \emph{sijo} and
        novels).
\end{itemize}

\section*{Summary}
\begin{itemize}
  \item Chinese characters are logographic; most are phono-semantic compounds organized
        by radicals.
  \item Literary Chinese was the shared written language of the region, like Latin.
  \item Korea and Japan adapted characters to languages with very different grammar,
        creating kana (by simplifying characters) and hangul (an invented alphabet).
  \item Shared Sino-vocabulary is a visible trace of the common heritage, much of it
        coined in Meiji Japan.
  \item Simplified vs.\ traditional characters is a twentieth-century political divide.
\end{itemize}
